\begin{lemma}[Theorem 2 in the nondegenerate case]
Let $n\ge2$, $p>1$ and $r>0$, suppose the Sobolev--Poincar\'e inequality on
balls holds in dimension $n$ (\noderef{CitedSobolevPoincareBall}), and assume
the cited regularity result (C1) (\noderef{CitedRegularityEpsPos}) for
\emph{every} field $A'\colon\R^n\to\R^n$, \emph{every} open set
$\Omega'\subseteq\R^n$ and all parameters $\varepsilon',\gamma_0',\gamma_1'$
with $\gamma_0'>0$ and $\gamma_1'>0$.  Then there is a constant
$\alpha=\alpha(p,r,n)$ with $0<\alpha\le1$ with the following property.  Let
$\varepsilon>0$, $\gamma_0>0$, $\gamma_1\le r\gamma_0$, let
$\Omega_0\subseteq\R^n$ be open, let $A$ satisfy (Q1) of the thesis--(Q2) of the thesis
on $\Omega_0$ with parameters $p,\varepsilon,\gamma_0,\gamma_1$
(\noderef{EllipticStructureQ12}), let $u$ be a weak solution of
$\operatorname{div}A(\nabla u)=0$ in $\Omega_0$ (\noderef{IsWeakSolutionOn})
with weak gradient $g$, let $x_0\in\R^n$ and $R_0>0$ with
$B_{R_0}(x_0)\subseteq\Omega_0$, and let $T\ge0$ satisfy $\|g\|\le T$ a.e.\ on
$B_{R_0}(x_0)$.  Then for every $0<\rho\le R_0$ and every $k$,
\[
\operatorname*{osc}\bigl(g^k,B_\rho(x_0)\bigr)
\ \le\ 16\,T\,\Bigl(\frac{\rho}{R_0}\Bigr)^{\alpha},
\]
where $g^k=D_ku$ is the $k$-th component of $g$ and
$\operatorname*{osc}(f,E)=\operatorname{ess\,sup}_Ef-
\operatorname{ess\,inf}_Ef$ (\noderef{EssOscOn}).

The hypothesis is the \emph{open} ball $B_{R_0}(x_0)\subseteq\Omega_0$ and not
$\overline B_{R_0}(x_0)\subseteq\Omega_0$: every node used in the proof is
applied on an open ball contained in $B_{R_0}(x_0)$, so nothing is lost, and the
weaker hypothesis is what lets a consumer apply the lemma with $\Omega_0$ itself
an open ball --- which is what the removal of the hypothesis $\varepsilon>0$
needs, since the approximating solutions it produces live only on a ball.

Both the oscillation and the bound $T$ are essential, as on p.~6 of the paper:
$g$ is a weak gradient, hence an almost-everywhere class, so no pointwise
assertion may be made about it.  Quantifying over all essential upper bounds $T$
is equivalent to using $\esssup_{B_{R_0}}\|g\|$, which is itself such a bound
whenever it is finite.  The constant $\alpha$ is bound before
$\varepsilon,\gamma_0,\gamma_1,T,\Omega_0,A,u,g,x_0,R_0,\rho,k$, and the
structure constants enter only through $\gamma_1\le r\gamma_0$: this is the
formal content of ``$\alpha$ depends only on $p$, $\gamma_1/\gamma_0$ and $n$''.
The factor $16$ is absolute.  Note that $\gamma_1>0$ is \emph{not} assumed; it
is derived from $\gamma_0>0$ and the structure conditions in (S0) below.

The hypothesis \noderef{CitedSobolevPoincareBall} is the cited result (C6) of
the paper's Note for the formalization, carried here as an explicit argument
because the proof uses it --- it enters through the constants of
\noderef{OscillationDecayLemma} and \noderef{DegenerateAlternative} --- which is
what the Note permits for the results the thesis quotes without proof.  Although
it is bound before the existential for $\alpha$, it introduces no dependence of
$\alpha$ on any quantity outside the list the paper names: (C6) is a statement
about $n$ alone, and $n$ is already bound first, so the hypothesis carries no
numerical content of its own.

The cited regularity result (C1) (\noderef{CitedRegularityEpsPos}) is assumed
for \emph{every} field $A'$, \emph{every} open set $\Omega'$ and all parameters
$\varepsilon',\gamma_0',\gamma_1'$ with $\gamma_0'>0$ and $\gamma_1'>0$.  That
is the form (C1) has in the paper's Note, which states it for an arbitrary
operator and an arbitrary solution: ``If $\varepsilon>0$ and
$A\colon\R^n\to\R^n$ satisfies (Q1) of the thesis and (Q2) of the thesis, then every weak
solution $u\in W^{1,p}_{\mathrm{loc}}(\Omega)$ of (3.1) of the thesis satisfies
$u\in W^{2,2}_{\mathrm{loc}}(\Omega)$ and $\nabla u\in C(\Omega)$'' (Note for
the formalization, (C1)).  So the universal clause is the faithful
transcription, and a clause attached only to the field $A$ and the domain
$\Omega_0$ of this lemma would be the weakening.  The two positivity side
conditions and the openness condition are not optional: the structure conditions
(Q1) of the thesis--(Q2) of the thesis carry positive constants in the thesis
(``$\gamma_0,\gamma_1>0$''), and (C1) is a statement about weak solutions on an
open set, so a clause quantified over $\gamma_0'\le0$ or over arbitrary
$\Omega'$ would assert more than the cited result does --- and since this is a
hypothesis of the present lemma, asserting more than the literature gives risks
making the lemma vacuously true.  The remaining antecedents need no guard,
because \noderef{CitedRegularityEpsPos} carries $\varepsilon>0$,
(Q1) of the thesis--(Q2) of the thesis and the weak-solution property inside its own
statement.  Quantifying the set is not decorative: the proof consumes (C1) on
the balls $B_{R'}\subseteq B_{R_0}$ on which the alternative
\noderef{DegenerateAlternative} is applied --- that node asks for the predicate
itself at the ball --- and \noderef{CaccioppoliLevelWeight} takes the family
$D(\nabla u)$ as a parameter, which the proof must produce from the existential
inside (C1)'s conclusion.  Like (C6), (C1) carries no numerical content: it
asserts membership and the existence of a continuous representative of the
gradient, and names no constant, so although it is bound before the existential
for $\alpha$ it introduces no dependence of $\alpha$ on anything outside the
list $p$, $\gamma_1/\gamma_0$, $n$ that the paper names.
\end{lemma}
\begin{proof}
Write $B_\sigma=B_\sigma(x_0)$, $w=\varepsilon+|\nabla u|^2$ and $g^k=D_ku$.
Constants named $c$, $C$ or $\gamma$ with a subscript are positive and depend
only on $n$, $p$ and $r$; this is checked at each occurrence.
$\omega_n=|B_1|$, so $|B_\sigma|=\omega_n\sigma^{\,n}$ (Preamble: the volume of
a ball).  This is thesis pp.~27--29.

\medskip
\noindent\emph{Step 0: the constants, chosen before all the data.}

\textbf{(S0) A generic positivity claim for the second structure constant.}
\emph{Claim.  Let $\Omega_\ast\subseteq\R^n$ be open with $|\Omega_\ast|>0$, let
$\varepsilon_\ast\ge0$ and $\gamma_{0\ast},\gamma_{1\ast}\in\R$ with
$\gamma_{0\ast}>0$, and let $A_\ast$ satisfy (Q1) of the thesis--(Q2) of the thesis on
$\Omega_\ast$ with parameters $p,\varepsilon_\ast,\gamma_{0\ast},\gamma_{1\ast}$
(\noderef{EllipticStructureQ12}).  Then $\gamma_{0\ast}\le\gamma_{1\ast}$, and in
particular $\gamma_{1\ast}>0$.}

The claim is a self-contained consequence of
(Q1) of the thesis--(Q2) of the thesis alone and is stated for an arbitrary set, field
and parameter triple rather than for the $\Omega_0,A,\varepsilon,\gamma_0,
\gamma_1$ of the present lemma; exactly one instance is taken below, in (B0),
at the lemma's own data, and that one instance is then quoted at three later
places.  The generic form is recorded only so that the instance is visibly
independent of everything else in the proof.

\emph{Proof of the claim.}  By \noderef{EllipticStructureQ12} the conditions
(Q1) of the thesis and (Q2) of the thesis hold for almost every $x\in\Omega_\ast$.  Their
two exceptional sets are null, so their union is null, while
$|\Omega_\ast|>0$; hence the union cannot contain the whole of $\Omega_\ast$ and
there is a single $x\in\Omega_\ast$ at which both hold.  Fix such an $x$, take
$z=0$, $h=e_1$ and $\xi=e_1$.  Then
$\varepsilon_\ast+|h|^2=\varepsilon_\ast+1\ne0$, which is the side condition
under which (Q1) of the thesis and (Q2) of the thesis are asserted, and $|\xi|^2=1$, so
they give
\[
\gamma_{0\ast}(\varepsilon_\ast+1)^{\frac p2-1}
\le\bigl(D_hA_\ast(x,0,h)e_1,e_1\bigr)=\bigl(D_hA_\ast(x,0,h)e_1\bigr)^1
\le\bigl|\bigl(D_hA_\ast(x,0,h)e_1\bigr)^1\bigr|
\le\gamma_{1\ast}(\varepsilon_\ast+1)^{\frac p2-1},
\]
the middle equality because $\xi=e_1$ has only the first coordinate nonzero, and
the last inequality being (Q2) of the thesis at $i=j=1$, whose $j$-th coordinate
vector is exactly the $\xi=e_1$ used in (Q1) of the thesis.  Now
$\varepsilon_\ast+1>0$, so $(\varepsilon_\ast+1)^{p/2-1}>0$, and dividing gives
$\gamma_{0\ast}\le\gamma_{1\ast}$; with $\gamma_{0\ast}>0$ this gives
$\gamma_{1\ast}>0$.  \emph{End of proof of the claim.}

The claim is needed because the carried clause (C1)
(\noderef{CitedRegularityEpsPos}) is assumed only for positive structure
constants --- its antecedents are $\gamma_0'>0$ and $\gamma_1'>0$ --- while the
data available at each consuming site give only $\gamma_{0\ast}>0$ and a ratio
bound $\gamma_{1\ast}\le r\gamma_{0\ast}$, which permits $\gamma_{1\ast}\le0$.

Note the shape of the instance: it is taken at $h=e_1$, \emph{not} at $h=0$.  The
weight is $(\varepsilon_\ast+1)^{p/2-1}$, which is positive for every
$\varepsilon_\ast\ge0$, so the claim covers $\varepsilon_\ast=0$ as well as
$\varepsilon_\ast>0$ --- which is why it is stated with $\varepsilon_\ast\ge0$,
even though the only instances taken below have $\varepsilon_\ast>0$.  At $h=0$
and $\varepsilon_\ast=0$ the weight would be $0^{p/2-1}$, which the run's rpow
convention makes $0$ for $p\ne2$, and the two conditions would read
$0\le\cdot\le0$ and yield nothing.

\textbf{(S1)} Apply \noderef{DegenerateAlternative} with $n,p,r$ and the
Sobolev--Poincar\'e hypothesis \noderef{CitedSobolevPoincareBall} of the present
lemma, and let $\theta_2,\nu\in(0,1)$ be the constants it produces.  Put
\[
\alpha_1=\min\Bigl\{\frac{\log(1/\nu)}{\log2},\,1\Bigr\}\in(0,1],
\qquad\text{so}\qquad \nu\le2^{-\alpha_1},
\]
the implication because $\alpha_1\log2\le\log(1/\nu)$.

\textbf{(S2)} Apply \noderef{CaccioppoliLevelWeight} with $n,p,r$ and let
$c_4>0$ be its constant.  Put
\[
c_p=\min\{\tfrac1{16},\,2^{-1-4/p}\}\in(0,\tfrac1{16}],
\qquad
K_p=\frac{\tfrac14+n}{c_p^{\,2}} .
\]
Apply \noderef{PinchedCaccioppoliLevelBound} with $n$, the exponent
$s=p/2-1$, the ratio $K=K_p$ and the Caccioppoli constant $c=c_4$ --- its
hypotheses are $1\le n$, $K_p\ge1$ and $c_4>0$ --- and let
\[
\gamma_\sharp>0
\]
be the constant $\gamma$ it produces; it depends only on $n$, $s=p/2-1$, $K_p$
and $c_4$, hence only on $n,p,r$.
Here $c_p$ is the constant in part (a) of \noderef{DegenerateAlternative}, and
$K_p\ge1+4n$ because $4c_p^{\,2}\le4/256<1$ and
$K_p=(1+4n)/(4c_p^{\,2})$; this is why one $\gamma_\sharp$ covers both pinches
below.
\emph{Note that $K_p$ and not $64(1+4n)$ is the right ratio}: $64(1+4n)$ is
$K_p$ evaluated at $c_p=1/16$, which is the paper's constant in part (a) before
the correction reported in \noderef{DegenerateAlternative}, and it is correct
exactly for $p\ge4/3$; for $1<p<4/3$ one has $c_p=2^{-1-4/p}<1/16$ and $K_p$ is
strictly larger.  $K_p$ is still a function of $n$ and $p$ alone, so no
dependence list moves.

\textbf{(S3)} Apply \noderef{OscillationDecayLemma} with $n$, the constant
$\gamma_\sharp$, the ratio $\lambda_1=\tfrac23\in(0,1)$ and
\noderef{CitedSobolevPoincareBall}, and let $\delta\in(0,1)$ be the constant it
produces.  Put
\[
\alpha_2=\min\Bigl\{\frac{\log(1/\delta)}{\log2},\,1\Bigr\}\in(0,1],
\qquad
\alpha=\min\{\alpha_1,\alpha_2\}\in(0,1] .
\]

All of $\theta_2,\nu,\alpha_1,c_4,c_p,K_p,\gamma_\sharp,\delta,\alpha_2,\alpha$
are functions of $n$, $p$ and $r$ alone and are fixed before
$\varepsilon,\gamma_0,\gamma_1,T,\Omega_0,A,u,g,x_0,R_0,\rho,k$, which is the
quantifier order the statement asserts.  The order of choice is acyclic:
$\theta_2,\nu,c_4$ first, then $\alpha_1,c_p,K_p,\gamma_\sharp$, then $\delta$,
then $\alpha_2$ and $\alpha=\min\{\alpha_1,\alpha_2\}$.

Now fix data as in the statement.

\medskip
\noindent\emph{Step 1: the estimate.}

\textbf{(B0r) Restriction to a ball.}  We record once two passages used at
several places below.

\emph{The weak-solution property restricts.}  Let $\Omega_\ast$ be open, let
$B=B_\sigma(x_0)\subseteq\Omega_\ast$ with $\sigma>0$, and suppose $u$ is a weak
solution of $\operatorname{div}A(\nabla u)=0$ in $\Omega_\ast$
(\noderef{IsWeakSolutionOn}) with weak gradient $g$.  Then $u$ is a weak
solution of the same equation in $B$ with weak gradient $g$.  This is
\noderef{IsWeakSolutionRestrict}, applied with $n$, the exponent $p$, the open
set $\Omega_\ast$, the centre $x_0$, the radius $\sigma$, the first field in the
form $(x,z,h)\mapsto A(h)$, the second field identically zero, and the pair
$(u,g)$; its hypotheses are $1\le n$ and $1\le p$ (from $2\le n$ and $1<p$), the
openness of $\Omega_\ast$, $0<\sigma$, the inclusion $B\subseteq\Omega_\ast$ and
the weak-solution property on $\Omega_\ast$, all of which are in hand at each
use below.

\emph{Likewise (Q1) of the thesis--(Q2) of the thesis restrict}: let $\Omega_\ast$ be any
set and $B\subseteq\Omega_\ast$.  Each of the \emph{last four} clauses of
\noderef{EllipticStructureQ12} is an
assertion of the form ``for (almost) every $x$, if $x\in\Omega_\ast$ then
\dots'', and replacing $\Omega_\ast$ by a subset only weakens the antecedent;
the first two, the Carath\'eodory pair, carry no domain guard --- measurability
of $x\mapsto A(x,z,h)$ is asserted on $\R^n$ and joint continuity of
$(z,h)\mapsto A(x,z,h)$ at every $x\in\R^n$ --- so they are the same assertion
for $B$ as for $\Omega_\ast$ and need no restriction at all.  All six therefore
hold for $B$.

\textbf{(B0) Data on the ball, and the maximal partial derivative.}  By (B0r)
with $\Omega_\ast=\Omega_0$ and $B=B_{R_0}$, together with clause (2) of
\noderef{WeakGradientRestrict}, the structure
(Q1) of the thesis--(Q2) of the thesis, the weak-solution property and
$u\in W^{1,p}_{\mathrm{loc}}$ are all available on $B_{R_0}$.  Applying
the present lemma's universal hypothesis (C1)
(\noderef{CitedRegularityEpsPos}) at $\varepsilon>0$, $\gamma_0,\gamma_1$ and
the open set $B_{R_0}$ --- legitimate because that clause requires
$\gamma_0'>0$, which is a datum of the present lemma, and
$\gamma_1'>0$, which is (S0) applied to the open set $\Omega_\ast=\Omega_0$,
the field $A_\ast=A$ and the parameters
$\varepsilon_\ast=\varepsilon,\gamma_{0\ast}=\gamma_0,\gamma_{1\ast}=\gamma_1$
of the present lemma; its hypotheses hold because $\Omega_0$ is open and
contains $B_{R_0}$ with $R_0>0$, so $|\Omega_0|\ge\omega_nR_0^{\,n}>0$
(Preamble: the volume of a ball), because $\gamma_0>0$, and because $A$
satisfies (Q1) of the thesis--(Q2) of the thesis on $\Omega_0$ ---
and using the structure and the weak-solution property
just obtained, we get: a field $g'$ with $g'$ continuous on $B_{R_0}$ and
$g=g'$ almost everywhere on $B_{R_0}$, together with a family
$Dg=(Dg^i)_i$ with $g'^i\in W^{1,2}_{\mathrm{loc}}(B_{R_0})$ of weak gradient
$Dg^i$ for every $i$ (\noderef{MemW1pLocOn}).  \emph{The continuity is of the
representative $g'$, not of the given $g$}: a weak gradient is determined only
almost everywhere, so continuity is a property of a representative, and (C1)
says so.  The $W^{1,2}_{\mathrm{loc}}$ clause is stated for $g'$ and transfers
to $g$ unchanged, $\mathrm{MemW1pLocOn}$ being invariant under modification on a
null set (clause (3) of \noderef{AEModificationInvariance}, applied with $q=2$,
the open set $B_{R_0}$, the pair $(g'^i,g^i)$ and the field $Dg^i$, whose
hypotheses $1\le n$, $1\le2$, the openness of $B_{R_0}$ and the
almost-everywhere equality are all in hand).  Fix this family once; by clause
(2) of \noderef{WeakGradientRestrict} it serves every open subball as well.

For $0<\sigma\le R_0$ put
\[
\mu(\sigma)=\max_{1\le k\le n}\operatorname*{ess\,sup}_{B_\sigma}|g^k| .
\]
The maximum is over the finitely many $k$, and each inner quantity is an
essential supremum, which is the paper's own reading: l.~290--291 declares that
``Max'' in (2.1) of the thesis actually means ``essential max'', a convention the
paper follows from then on, and l.~293--294 accordingly sets
$M=\operatorname*{ess\,sup}_{B_R}f$.  Three properties of $\mu$ are used below, and all three come
from the definition of the essential supremum, exactly as (O) below derives the
monotonicity of $\operatorname*{osc}$ from the definition of \noderef{EssOscOn}.

First, \emph{$\mu(\sigma)$ is a finite nonnegative real number}, so that none of
the quantities below is a junk value.  Each ball $B_\sigma$ with $0<\sigma\le R_0$
has volume $\omega_n\sigma^n>0$ (Preamble: the volume of a ball).  And
$|g^k|\le\|g\|\le T$ a.e.\ on $B_{R_0}$, the first inequality pointwise because
$|g^k|$ is one coordinate of the vector $g$; since
$B_\sigma\subseteq B_{R_0}$, the volume measure restricted to $B_\sigma$ is
dominated by the one restricted to $B_{R_0}$, so a null set for the latter is null
for the former and the bound holds a.e.\ on $B_\sigma$ as well.  Hence each
$|g^k|$ is essentially bounded above on $B_\sigma$ and its essential supremum
there is a real number at most $T$; it is also $\ge0$, because $|g^k|\ge0$ and an
essential upper bound on a set of positive measure must be $\ge0$ --- the set
where $|g^k|$ does not exceed it is of positive measure, hence nonempty.  The
maximum of the $n$ resulting real numbers is $\mu(\sigma)$, so
$0\le\mu(\sigma)\le T$.

Second, \emph{$|g^k|\le\mu(\sigma)$ a.e.\ on $B_\sigma$ for every $k$.}  For each
$k$ the essential supremum is an essential upper bound, i.e.\
$|g^k|\le\operatorname*{ess\,sup}_{B_\sigma}|g^k|$ a.e.\ on $B_\sigma$, and
$\operatorname*{ess\,sup}_{B_\sigma}|g^k|\le\mu(\sigma)$ because $\mu(\sigma)$ is
the maximum over $k$.  The union of the $n$ exceptional null sets is null, so
off that union the bound holds for all $k$ at once; hence it holds a.e.\ on
$B_\sigma$ simultaneously for all $k$.

Third, \emph{$\mu(\sigma)$ is the least such number} --- which is what
\noderef{DegenerateAlternative} requires of its $\mu$.  Let $c\in\R$ satisfy
$|g^k|\le c$ a.e.\ on $B_\sigma$ for every $k$.  Then for each $k$ the number $c$
is an essential upper bound of $|g^k|$ on $B_\sigma$, and the essential supremum
is by definition the \emph{least} essential upper bound, so
$\operatorname*{ess\,sup}_{B_\sigma}|g^k|\le c$.  Taking the maximum over the
finitely many $k$ gives $\mu(\sigma)\le c$.

Finally, \emph{$\mu$ is nondecreasing}, and this now follows from the previous
two.  Let $0<\sigma\le\sigma'\le R_0$, so $B_\sigma\subseteq B_{\sigma'}$.  By the
second property at radius $\sigma'$, $|g^k|\le\mu(\sigma')$ a.e.\ on
$B_{\sigma'}$ for every $k$, hence --- by the same domination of restricted
measures used above --- a.e.\ on $B_\sigma$ for every $k$.  So $\mu(\sigma')$ is
a number of the kind the third property quantifies over at radius $\sigma$, and
that property gives $\mu(\sigma)\le\mu(\sigma')$.  In particular
$\mu(\sigma)\le\mu(R_0)\le T$, the last inequality being the first property at
$\sigma=R_0$.  No case distinction on $T$ is needed anywhere below: every
estimate is run on $\mu$ and converted at the very end through
$0\le\mu(R_0)\le T$, which holds for $T=0$ as well.

\textbf{(B1) The dyadic scales and the alternative.}  Put $R_m=2^{-m}R_0$ and
$\mu_m=\mu(R_m)$ for $m=0,1,2,\dots$.  For each $m$ apply
\noderef{DegenerateAlternative} --- the instance of (S1) --- with the data
$\varepsilon,\gamma_0,\gamma_1$, $\mu_m$, the field $A$, the functions $u,g$,
the family $Dg$, the centre $x_0$ and the radius $R_m$.  Its hypotheses are
$0<\varepsilon$, $0<\gamma_0$, $\gamma_1\le r\gamma_0$, $0<R_m$,
(Q1) of the thesis--(Q2) of the thesis on $B_{R_m}$, the cited regularity (C1) on
$B_{R_m}$ --- the present lemma's universal hypothesis at that open set,
instantiated at $\gamma_0>0$ and at $\gamma_1>0$, the latter by the instance of
(S0) taken in (B0) ---
the weak-solution property on $B_{R_m}$ --- by (B0r) with
$\Omega_\ast=B_{R_0}$ and $B=B_{R_m}$ --- the family
$g^i\in W^{1,2}_{\mathrm{loc}}(B_{R_m})$, by clause (2) of
\noderef{WeakGradientRestrict} with $q=2$, the bound $|g^k|\le\mu_m$ a.e.\ on
$B_{R_m}$ for every $k$ and its minimality.  Write
\[
H_m\ :\ \exists k,\ \ \theta_2|B_{R_m}|\le\bigl|B_{R_m}\cap\{g^k>\tfrac{\mu_m}2\}\bigr|
\ \text{ or }\ \theta_2|B_{R_m}|\le\bigl|B_{R_m}\cap\{g^k<-\tfrac{\mu_m}2\}\bigr| .
\]
If $H_m$ fails then part (b) gives $|g^k|\le\nu\mu_m$ a.e.\ on $B_{R_m/2}
=B_{R_{m+1}}$ for every $k$, hence
\[
\mu_{m+1}\le\nu\,\mu_m .
\tag{3.25}
\]
If $H_m$ holds, with witness $k_0$, then part (a) at that $k_0$ gives
\[
|g^{k_0}|\ \ge\ c_p\,\mu_m\qquad\text{a.e.\ on }B_{R_{m+1}} .
\tag{3.26}
\]
(The paper states (3.26) of the thesis for $|\nabla u|$; the bridge is the
componentwise inequality $|g^{k_0}|\le\|\nabla u\|$, so the displayed bound for
the single maximising component implies the paper's bound for $\|\nabla u\|$,
and it is the component form that part (a) delivers.)

Two cases.

\textbf{(B2) Case I: $H_m$ fails for every $m$.}  Then (3.25) holds for all $m$,
so $\mu_m\le\nu^m\mu_0\le2^{-\alpha_1m}\mu(R_0)$ by (S1).  Given
$0<\rho\le R_0$ choose the integer $m\ge0$ with $R_{m+1}<\rho\le R_m$ (take
$m=0$ if $\rho=R_0$; such an $m$ exists because $R_m\downarrow0$).  Then
$2^{-m}=R_m/R_0<2\rho/R_0$, so
\[
\mu(\rho)\le\mu_m\le2^{-\alpha_1m}\mu(R_0)
=\Bigl(\frac{R_m}{R_0}\Bigr)^{\alpha_1}\mu(R_0)
\le\Bigl(\frac{2\rho}{R_0}\Bigr)^{\alpha_1}\mu(R_0)
\le2\Bigl(\frac{\rho}{R_0}\Bigr)^{\alpha}\mu(R_0),
\tag{B.1}
\]
using $\alpha_1\le1$ for $2^{\alpha_1}\le2$ and $\alpha\le\alpha_1$ together
with $\rho/R_0\le1$ for $(\rho/R_0)^{\alpha_1}\le(\rho/R_0)^{\alpha}$.  Since
$\operatorname*{osc}(g^k,B_\rho)\le2\mu(\rho)$ for every $k$ --- both
$\operatorname{ess\,sup}$ and $-\operatorname{ess\,inf}$ of $g^k$ on $B_\rho$
are at most $\mu(\rho)$ --- this gives
$\operatorname*{osc}(g^k,B_\rho)\le4T(\rho/R_0)^{\alpha}$, which is
the asserted bound.

\textbf{(B3) Case II: the first $m_0$ with $H_{m_0}$.}  Let $m_0$ be least with
$H_{m_0}$ true.  Then (3.25) holds for $m=0,\dots,m_0-1$ and (3.26) holds at
$m_0$ with some witness $k_0$.  In particular
\[
\mu_{m}\le\nu^{m}\mu(R_0)\le2^{-\alpha_1m}\mu(R_0)\qquad(0\le m\le m_0),
\tag{B.2}
\]
and $\mu_{m_0}>0$: if $\mu_{m_0}=0$ then $g^k=0$ a.e.\ on $B_{R_{m_0}}$ for
every $k$, so both sets in $H_{m_0}$ are null while
$\theta_2|B_{R_{m_0}}|>0$, contradicting $H_{m_0}$.  Put
\[
\mathcal B=B_{R_{m_0+1}},\qquad \mu=\mu_{m_0}>0,\qquad \lambda=\frac\mu2>0 .
\]

\textbf{(B4) The pinch on $\mathcal B$.}  On $\mathcal B$ we have, a.e.,
$|g^k|\le\mu_{m_0+1}\le\mu$ for every $k$, hence
$\|\nabla u\|^2=\sum_k(g^k)^2\le n\mu^2$ and so
\[
\varepsilon\ \le\ w\ \le\ \varepsilon+n\mu^2
\qquad\text{a.e.\ on }\mathcal B .
\tag{B.3}
\]
Two sub-cases, and in both $w$ is pinched almost everywhere on \emph{all} of
$\mathcal B$ between two positive constants whose ratio is at most $K_p$.  (All
the bounds here and in the block below hold almost everywhere, which is all the
integral inequalities of the block need; the companion instantiation in
\noderef{DegenerateAlternative} has them everywhere on the ball because (C1)
supplies a continuous \emph{representative} $g'$ of the gradient there, but that
is not used here.)
\begin{itemize}
\item If $\varepsilon\ge\lambda^2=\mu^2/4$, then by (B.3)
  $\varepsilon\le w\le\varepsilon+4n\varepsilon=(1+4n)\varepsilon$, so
  $\kappa_1=\varepsilon>0$, $\kappa_2=(1+4n)\varepsilon$ and $K=1+4n\le K_p$.
\item If $\varepsilon<\lambda^2=\mu^2/4$, then by (3.26) and
  $|g^{k_0}|\le\|\nabla u\|$ we have $w\ge\|\nabla u\|^2\ge c_p^2\mu^2$ a.e.\ on
  $\mathcal B$, while (B.3) gives $w\le\mu^2/4+n\mu^2=(\tfrac14+n)\mu^2$.  So
  $\kappa_1=c_p^2\mu^2>0$, $\kappa_2=(\tfrac14+n)\mu^2$ and
  $K=(\tfrac14+n)/c_p^2=K_p$.
\end{itemize}
This is the step at l.~944--945 of the paper, ``Since we also have
$|\nabla u|\le\sqrt n\,\mu(R_{m_0})$ in $B_{R_{m_0+1}}$ it follows from
(3.4) of the thesis that (3.9) of the thesis holds in $B_{R_{m_0+1}}$'', with the
correction of the constant: the lower pinch is $c_p^2\mu^2$, not
$(1/16)^2\mu^2=\mu^2/256$, because part (a) of
\noderef{DegenerateAlternative} delivers $c_p\mu$ and not $\mu/16$.

\textbf{(B5) The Caccioppoli inequality and the instances of
\textup{(2.1)}.}  Apply \noderef{CaccioppoliLevelWeight} --- the instance of
(S2), at $n,p,r$ with the constant $c_4$.  Its remaining hypotheses are
$0<\varepsilon$, $0<\gamma_0$, $\gamma_1\le r\gamma_0$, the openness of
$\mathcal B$, (Q1) of the thesis--(Q2) of the thesis (\noderef{EllipticStructureQ12}), the
cited regularity (C1) (\noderef{CitedRegularityEpsPos}) --- the universal clause
at the open set $\mathcal B$, whose side conditions $\gamma_0'>0$ and
$\gamma_1'>0$ are the datum and the instance of (S0) taken in (B0) --- the weak
solution
property (\noderef{IsWeakSolutionOn}) and
$g^i\in W^{1,2}_{\mathrm{loc}}(\mathcal B)$ for every $i$
(\noderef{MemW1pLocOn}), all available on $\mathcal B\subseteq B_{R_0}$ by (B0)
and (B0r).  Hence for every $k$, every $\ell\in\R$ and every $\zeta$ smooth with
compact support in $\mathcal B$,
\[
\int_{\mathcal B\cap\{g^k>\ell\}}w^{\frac p2-1}\|Dg^k\|^2\zeta^2\,dx
\ \le\ c_4\int_{\mathcal B}w^{\frac p2-1}\bigl((g^k-\ell)^+\bigr)^2|\nabla\zeta|^2\,dx
\tag{3.4$^+$}
\]
and
\[
\int_{\mathcal B\cap\{g^k<\ell\}}w^{\frac p2-1}\|Dg^k\|^2\zeta^2\,dx
\ \le\ c_4\int_{\mathcal B}w^{\frac p2-1}\bigl((\ell-g^k)^+\bigr)^2|\nabla\zeta|^2\,dx .
\tag{3.4$^-$}
\]
These are the two conclusions of that node verbatim, with
$\sum_i(\partial_i\zeta)^2=|\nabla\zeta|^2$ and
$\max(g^k-\ell,0)=(g^k-\ell)^+$.

Now apply \noderef{PinchedCaccioppoliLevelBound} --- the instance of (S2), with
the constant $\gamma_\sharp$ --- at the centre $x_0$, the radius $R_{m_0+1}$ (so
that its ball is $\mathcal B$), the bound $\Theta=\mu$, a level $\ell$, the pinch
constants $\kappa_1,\kappa_2$ of (B4), the function $f=g^k$, the field
$G=Dg^k$ and a weight $\widetilde w$ fixed as follows.  That node asks for a
weight measurable on all of $\R^n$, while (B0) gives a \emph{representative}
$g'$ of the gradient, continuous only on $B_{R_0}$, with $g=g'$ almost
everywhere there --- (C1) asserts continuity of a representative and not of the
given $g$, a weak gradient being an almost-everywhere class.  So build the
weight from $g'$ and not from $g$:
\[
\widetilde w(x)=\varepsilon+|g'(x)|^2\ \text{ for }x\in\mathcal B,
\qquad
\widetilde w(x)=\varepsilon\ \text{ for }x\notin\mathcal B .
\]
Then $\widetilde w$ is measurable on $\R^n$ because
$\mathcal B$ is open, so $\{\mathcal B,\R^n\setminus\mathcal B\}$ is a
measurable partition on each piece of which $\widetilde w$ is measurable --- on
$\mathcal B\subseteq B_{R_0}$ because $g'$ is continuous there, and on the
complement because $\widetilde w$ is constant --- and $\widetilde w=w$ almost
everywhere on $\mathcal B$, which is all that is used, as the next sentence
records.  Every hypothesis and conclusion of that node reads its weight
only inside its ball, which here is $\mathcal B$: the two pinch hypotheses are
almost-everywhere statements on subsets of $\mathcal B$, the two Caccioppoli
hypotheses integrate over $\mathcal B$ and subsets of it, and the conclusions do
not mention the weight; so (B4) and (3.4$^\pm$), which are statements about $w$
on $\mathcal B$, are statements about $\widetilde w$ verbatim.  Its hypotheses
are then: $0<R_{m_0+1}$; $0<\kappa_1$ and $\kappa_2\le K\kappa_1$ with
$K\le K_p$, which is (B4) --- and increasing $K$ to $K_p$ is legitimate, since
$\kappa_2\le K\kappa_1\le K_p\kappa_1$; $\widetilde w$ measurable, just shown;
$g^k\in W^{1,2}_{\mathrm{loc}}(\mathcal B)$ with weak gradient $Dg^k$
(\noderef{MemW1pLocOn}), by (B0) and clause (2) of
\noderef{WeakGradientRestrict}; and $|g^k|\le\mu$ a.e.\ on $\mathcal B$, which is
(B3)--(B4).  Because (B4) pinches $w$ on \emph{all} of $\mathcal B$, both of that
node's hypotheses (H3$^+$) and (H3$^-$) hold at every level $\ell\in\R$, and
(H4$^+$), (H4$^-$) are (3.4$^\pm$) above, read with $s=p/2-1$ and $c=c_4$.  Its
two implications therefore give, for every $k$, every $\ell\in\R$ and all
$0<\rho_0<\rho_1<R_{m_0+1}$,
\[
\mathrm{DeGiorgiLevelBound}\ n\ \gamma_\sharp\ g^k\ Dg^k\ x_0\ \ell\ \rho_0\ \rho_1
\quad\text{and}\quad
\mathrm{DeGiorgiLevelBound}\ n\ \gamma_\sharp\ (-g^k)\ (-Dg^k)\ x_0\ (-\ell)\
\rho_0\ \rho_1
\]
(\noderef{DeGiorgiLevelBound}).  This is the paper's (3.9) of the thesis on
$B_{R_{m_0+1}}$, in the form \S2 consumes it.

\textbf{(B6) The oscillation decay \textup{(3.27)}.}  Let $m\ge m_0+1$ and fix
$k$.  Put $T_m=\tfrac34R_m$, so that $T_m<R_m\le R_{m_0+1}$ and
$\lambda_1T_m=\tfrac23\cdot\tfrac34R_m=\tfrac12R_m=R_{m+1}$.  Apply
\noderef{OscillationDecayLemma} --- the instance of (S3), at $n$,
$\gamma_\sharp$ and $\lambda_1=\tfrac23$ --- with the centre $x_0$, the radius
$T_m$, the function $f=g^k$, the field $Dg^k$, and
$M=\operatorname{ess\,sup}_{B_{T_m}}g^k$,
$m^\flat=\operatorname{ess\,inf}_{B_{T_m}}g^k$.  Its hypotheses:
\begin{itemize}
\item $0<T_m$;
\item $g^k\in W^{1,2}(B_{T_m})$ with weak gradient $Dg^k$
  (\noderef{MemW1pOn}), by clause (4) of \noderef{WeakGradientRestrict} applied
  with $q=2$ and the open sets $B_{R_0}\supseteq B_{T_m}$, whose hypotheses are
  $g^k\in W^{1,2}_{\mathrm{loc}}(B_{R_0})$ from (B0),
  $\operatorname{closure}B_{T_m}=\overline B_{T_m}$ compact and
  $\overline B_{T_m}\subseteq B_{R_0}$, valid because $T_m<R_m\le R_0$;
\item $|g^k|\le\mu(R_0)$ a.e.\ on $B_{T_m}$, which is the required essential
  bound;
\item $M$ and $m^\flat$ are by definition the essential supremum and infimum;
\item its hypothesis (1): the family of (2.1) of the thesis predicate instances for
  $g^k$ at the levels $l\in[(m^\flat+M)/2,M]$ and the radii
  $\rho_0,\rho_1\in[\lambda_1T_m,T_m]=[R_{m+1},T_m]$ with $\rho_0<\rho_1$.  This
  is the first of the two instances (B5) supplies, which hold for every level
  $\ell\in\R$ and at any radii $0<\rho_0<\rho_1<R_{m_0+1}$ --- a range that
  covers $\rho_1\le T_m=\tfrac34R_m\le\tfrac34R_{m_0+1}<R_{m_0+1}$ --- with the
  constant $\gamma_\sharp$;
\item its hypothesis (2), which is \emph{not} a predicate but the bare display
  (2.11) of the thesis, for every $l\in[m^\flat,(m^\flat+M)/2]$ and the same radii.
  (B5) supplies the \emph{predicate}
  $\mathrm{DeGiorgiLevelBound}\ n\ \gamma_\sharp\ (-g^k)\ (-Dg^k)\ x_0\ (-l)\
  \rho_0\ \rho_1$ (\noderef{DeGiorgiLevelBound}), read at the level $l$ so that
  its own level argument is $-l$, and the conversion from that predicate to the
  display is the following, which uses only the integrand identity --- the
  predicate's clause (1) is simply discarded, the display having no such clause.
  Apply \noderef{WeakGradientPositivePart} with $q=2$, the level $c=-l$ and the
  open set $B_{T_m}$ --- of finite volume (Preamble: the volume of a ball) ---
  and the pair $(-g^k,-Dg^k)$.  That pair lies in $W^{1,2}(B_{T_m})$: the second
  bullet above gives $g^k\in W^{1,2}(B_{T_m})$ with weak gradient $Dg^k$
  (\noderef{MemW1pOn}), by clause (4) of \noderef{WeakGradientRestrict}, and
  clause (2) of \noderef{WeakGradientLinear} at $q=2$, the open set $B_{T_m}$,
  $a=-1$, $b=0$ and $(f_1,g_1)=(f_2,g_2)=(g^k,Dg^k)$ turns it into
  $-g^k\in W^{1,2}(B_{T_m})$ with weak gradient $-Dg^k$.  The ball $B_{T_m}$, and
  not $\mathcal B$, is the right set here: (B0) gives membership on $\mathcal B$
  only in the \emph{local} class $W^{1,2}_{\mathrm{loc}}(\mathcal B)$, whereas
  \noderef{WeakGradientPositivePart} asks for $W^{1,2}$ of the set itself, and
  $\overline B_{T_m}\subseteq B_{R_0}$ is what makes clause (4) available.  No
  generality is lost, because the radii at which the display is needed satisfy
  $\rho_0<\rho_1\le T_m$, so $B_{\rho_1}\subseteq B_{T_m}$.  The conclusion is
  $\bigl((-g^k)-(-l)\bigr)^+=(l-g^k)^+\in W^{1,2}(B_{T_m})$ with weak gradient
  $-Dg^k$ on $\{-l<-g^k\}$ and $0$ elsewhere.  A weak gradient is determined
  almost everywhere (Preamble: the fundamental lemma of the calculus of
  variations, second consequence), so the display's $\nabla(l-g^k)^+$ is that
  field a.e., and pointwise a.e.
  $|\nabla(l-g^k)^+|^2=\|Dg^k\|^2\chi_{\{g^k<l\}}$, using $\|-Dg^k\|=\|Dg^k\|$
  and $\{-l<-g^k\}=\{g^k<l\}$.  Hence the display's integral over all of
  $B_{\rho_0}$ equals the predicate's clause-(2) integral of $\|-Dg^k\|^2$ over
  $B_{\rho_0}\cap\{-g^k>-l\}$, the integrand vanishing off that set; and the two
  right-hand sides are the same expression, since
  $\bigl((-g^k)-(-l)\bigr)^+=(l-g^k)^+$ pointwise and
  $\{x\in B_{\rho_1}:-g^k(x)>-l\}=\{x\in B_{\rho_1}:g^k(x)<l\}$.  So clause (2)
  of the predicate \emph{is} the display, and hypothesis (2) follows.
\end{itemize}
Its conclusion is
$\operatorname*{osc}(g^k,B_{\lambda_1T_m})\le\delta\operatorname*{osc}(g^k,B_{T_m})$,
i.e.\ $\operatorname*{osc}(g^k,B_{R_{m+1}})\le\delta\operatorname*{osc}(g^k,B_{T_m})$.
We now record, once and for the two further uses below, the monotonicity of the
oscillation in the set.

\emph{(O) If $0<\sigma\le\sigma'\le R_0$ then
$\operatorname*{osc}(g^k,B_\sigma)\le\operatorname*{osc}(g^k,B_{\sigma'})$.}
By \noderef{EssOscOn} the oscillation on a set $E$ is
$\operatorname{ess\,sup}_Eg^k-\operatorname{ess\,inf}_Eg^k$, both taken with
respect to the volume measure restricted to $E$, so this is a statement about
those two quantities and is proved from the definition.  First, $B_\sigma$ and
$B_{\sigma'}$ have positive volume, $\omega_n\sigma^n$ and $\omega_n\sigma'^n$
(Preamble: the volume of a ball), and $|g^k|\le\mu(R_0)$ a.e.\ on $B_{R_0}$ by
(B0), so on each of the two balls $g^k$ is essentially bounded above and below
and both the essential supremum and the essential infimum are finite real
numbers --- neither is a junk value.  Next, $B_\sigma\subseteq B_{\sigma'}$, so
the volume measure restricted to $B_\sigma$ is dominated by the one restricted to
$B_{\sigma'}$; hence a set that is null for the latter is null for the former,
and therefore any property that holds almost everywhere on $B_{\sigma'}$ holds
almost everywhere on $B_\sigma$.  Put
$S'=\operatorname{ess\,sup}_{B_{\sigma'}}g^k$ and
$I'=\operatorname{ess\,inf}_{B_{\sigma'}}g^k$.  Then $g^k\le S'$ a.e.\ on
$B_{\sigma'}$, hence a.e.\ on $B_\sigma$, and the essential supremum over
$B_\sigma$ is the least number with that property, so
$\operatorname{ess\,sup}_{B_\sigma}g^k\le S'$.  Dually $g^k\ge I'$ a.e.\ on
$B_{\sigma'}$, hence a.e.\ on $B_\sigma$, and the essential infimum over
$B_\sigma$ is the greatest such number, so
$\operatorname{ess\,inf}_{B_\sigma}g^k\ge I'$.  Subtracting the second
inequality from the first gives (O).

With $T_m\le R_m$, so that $B_{T_m}\subseteq B_{R_m}$, (O) gives
\[
\operatorname*{osc}(g^k,B_{R_{m+1}})\ \le\ \delta\,\operatorname*{osc}(g^k,B_{R_m})
\qquad(m\ge m_0+1),
\tag{3.27}
\]
with the \emph{same} $\delta$: the passage to the radius $T_m<R_m$ costs nothing,
because the conclusion of \noderef{OscillationDecayLemma} is a ratio of
oscillations and $\operatorname*{osc}$ is monotone by (O), unlike the volume ratios
that the applications of Lemmas 2.3 and 2.4 inside
\noderef{DegenerateAlternative} have to absorb.

\textbf{(B7) The assembly.}  Fix $k$ and $0<\rho\le R_0$.  Recall
$\operatorname*{osc}(g^k,B_\sigma)\le2\mu(\sigma)$ for every $\sigma$.  Three
ranges.

\emph{(i) $R_{m_0}\le\rho\le R_0$.}  Exactly as in (B.1), using (B.2) in place
of (3.25) for all $m$: choose $m$ with $R_{m+1}<\rho\le R_m$; then $m\le m_0$,
because $\rho\le R_m$ and $\rho\ge R_{m_0}$ force $R_m\ge R_{m_0}$.  So
(B.2) applies and
\[
\mu(\rho)\le\mu_m\le2^{-\alpha_1m}\mu(R_0)\le2\Bigl(\frac\rho{R_0}\Bigr)^{\alpha}\mu(R_0),
\tag{3.28}
\]
whence $\operatorname*{osc}(g^k,B_\rho)\le4(\rho/R_0)^{\alpha}\mu(R_0)$.

\emph{(ii) $R_{m_0+1}\le\rho\le R_{m_0}$.}  Then
$\operatorname*{osc}(g^k,B_\rho)\le2\mu(R_{m_0})\le4(R_{m_0}/R_0)^{\alpha}\mu(R_0)$
by (3.28) at $\rho=R_{m_0}$, and $R_{m_0}=2R_{m_0+1}\le2\rho$, so
$(R_{m_0}/R_0)^{\alpha}\le2^{\alpha}(\rho/R_0)^{\alpha}\le2(\rho/R_0)^{\alpha}$
and
\[
\operatorname*{osc}(g^k,B_\rho)\ \le\ 8\Bigl(\frac\rho{R_0}\Bigr)^{\alpha}\mu(R_0).
\tag{3.30}
\]

\emph{(iii) $0<\rho\le R_{m_0+1}$.}  Choose the integer $m\ge0$ with
$R_{m+1}<\rho\le R_m$; it exists and is unique because $m\mapsto R_m=2^{-m}R_0$
is strictly decreasing with $R_0\ge R_{m_0+1}\ge\rho$ and $R_m\to0$.  Then
$m\ge m_0+1$, because $R_{m+1}<\rho\le R_{m_0+1}$ gives $R_{m+1}<R_{m_0+1}$ and
hence $m+1>m_0+1$ by that same strict monotonicity.  Iterating (3.27) from
$m_0+1$ to $m$,
\[
\operatorname*{osc}(g^k,B_\rho)\le\operatorname*{osc}(g^k,B_{R_m})
\le\delta^{\,m-m_0-1}\operatorname*{osc}(g^k,B_{R_{m_0+1}}),
\]
the first step by (O), the monotonicity of $\operatorname*{osc}$ established in
(B6), applied with $\sigma=\rho\le\sigma'=R_m$.  By
(S3) $\delta\le2^{-\alpha_2}$, so
$\delta^{\,m-m_0-1}\le2^{-\alpha_2(m-m_0-1)}=(R_m/R_{m_0+1})^{\alpha_2}$, and
$R_m<2\rho$ gives $(R_m/R_{m_0+1})^{\alpha_2}\le
2^{\alpha_2}(\rho/R_{m_0+1})^{\alpha_2}\le2(\rho/R_{m_0+1})^{\alpha_2}$.  With
$\operatorname*{osc}(g^k,B_{R_{m_0+1}})\le2\mu(R_{m_0+1})\le2\mu(R_{m_0})$ and
(3.28) at $\rho=R_{m_0}$,
\[
\operatorname*{osc}(g^k,B_\rho)
\le2\Bigl(\frac\rho{R_{m_0+1}}\Bigr)^{\alpha_2}\cdot2\cdot
2\Bigl(\frac{R_{m_0}}{R_0}\Bigr)^{\alpha}\mu(R_0)
=8\Bigl(\frac\rho{R_{m_0+1}}\Bigr)^{\alpha_2}
\Bigl(\frac{R_{m_0}}{R_0}\Bigr)^{\alpha}\mu(R_0).
\]
Since $\rho\le R_{m_0+1}$ and $\alpha\le\alpha_2$,
$(\rho/R_{m_0+1})^{\alpha_2}\le(\rho/R_{m_0+1})^{\alpha}
=(2\rho/R_{m_0})^{\alpha}\le2(\rho/R_{m_0})^{\alpha}$, so the product of the two
powers is at most $2(\rho/R_{m_0})^{\alpha}(R_{m_0}/R_0)^{\alpha}
=2(\rho/R_0)^{\alpha}$ and
\[
\operatorname*{osc}(g^k,B_\rho)\ \le\ 16\Bigl(\frac\rho{R_0}\Bigr)^{\alpha}\mu(R_0).
\tag{3.29}
\]
This is the paper's $\alpha=\min(\alpha_1,\alpha_2)$ assembly
(3.28) of the thesis--(3.30) of the thesis; the paper writes $C_1$ in the display at
l.~962 and $C_2$ in the one that follows, and writes
``$\lambda_1/\lambda_0$'' at l.~953 where it means $\gamma_1/\gamma_0$.  In all
three ranges $\operatorname*{osc}(g^k,B_\rho)\le16(\rho/R_0)^{\alpha}\mu(R_0)
\le16\,T\,(\rho/R_0)^{\alpha}$, which with Case I proves
the lemma.

\end{proof}
